\section*{Limitations}

\paragraph{Pretraining exposure and evaluation scope.} The LLMs we evaluate were pretrained on very large \LaTeX{} corpora and have seen essentially no \texttt{.tmu} during pretraining, placing \texttt{.tmu} at a substantial model-familiarity disadvantage. Evaluating models with better-matched exposure is future work. This scopes our evaluation, not our thesis: the architectural frictions {\LaTeX} poses for LLMs (one-off batch compilation, weak structural semantics, error messages detached from their root cause; Appendices~\ref{app:arch}--\ref{app:con}) are properties of the format, not artifacts of model familiarity, and remain our core motivation.

\paragraph{Evaluation scale and statistics.} The LLM-task evaluation is focused rather than broad: ${\sim}20$ items each for structure-locating and debugging, only a handful of merging tasks, drawn largely from a single source document (arXiv:2502.17655). We report point estimates without significance testing or variance, and do not yet claim domain generality. A larger, multi-domain benchmark with repeated trials and significance testing is future work.

\paragraph{Scope of the fine-tuning evidence.} The fine-tuning result uses a single LoRA run on one model family (Qwen2.5-7B-Instruct), a single seed, and formula-level rather than full-document data. Lower training loss is consistent with but does not by itself establish better downstream generation; validation loss, exact-match and rendered-equivalence metrics, and multi-seed runs across model sizes are future work.

\paragraph{Utility metric.} Our utility scores ($u_s$, $u_m$, $u_d$) combine task correctness with token cost under a hand-chosen 10k-token scale; absolute scores are not comparable across task types, and the metric entangles accuracy with token efficiency. Reporting these two axes separately is future work.

\paragraph{Baselines and the tool-vs-representation confound.} Our baselines are limited to {\LaTeX}; comparison against Typst, Markdown-with-math, and MathML/XML and other structured editors is future work. We note Markdown specifically because it reuses {\LaTeX} math syntax and is the natural lightweight baseline, but a fair comparison is non-trivial given the gap in intended use, so we do not treat it as central. A further confound is that some gains may stem from the editor's tooling (incremental update, structured editing, the {\LaTeX} import engine) rather than the representation alone; disentangling these requires controlled ablations we have not yet run.

\paragraph{On low entropy.} We currently support the low-entropy claim with an operational proxy (reduced source-encoding multiplicity for a fixed rendered output, e.g.\ \texttt{x\^{}2} vs.\ \texttt{x\^{}\{2\}}) together with indirect evidence from fine-tuning convergence. A full information-theoretic treatment is future work.

\paragraph{Future work.} A concrete agenda: (i) a fair controlled comparison with an ablation separating the contribution of tree structure from that of pre-expanded reference numbers, isolating the structure-locating input asymmetry (in \texttt{.tmu}, numbering is stored inline after incremental update, whereas {\LaTeX} requires a full scan); (ii) a larger multi-domain evaluation with repeated trials, significance testing, and additional baselines (Typst, Markdown-with-math, MathML); (iii) fine-tuning with validation loss, exact-match, and multi-seed runs across model sizes, plus a Markdown comparison and a larger, standardized end-to-end LLM-repair benchmark.

% \paragraph{Deployment evidence.} As an industry-track contribution, we stress that the real-world deployment evidence in this version is preliminary; broader production analysis of cost, latency, and reliability at scale is ongoing and not yet reported here.

% \paragraph{Human-facing evidence.} The \texttt{.tmu} lineage also supports WYSIWYG authoring, but this paper does not report controlled human-subject speed measurements or production deployment statistics. We therefore treat human editing benefits as motivation and architectural context, not as an empirical contribution.
